\documentclass[a4paper,12pt]{article}
\usepackage[utf8]{inputenc}
\usepackage[T1]{fontenc}
\usepackage[french]{babel}
\usepackage{amsmath,amsfonts,amsthm,amssymb,amsbsy,amscd,mathrsfs}
\usepackage{lmodern}
\usepackage{fourier}
%\usepackage{times,mathptmx}
%\DeclareSymbolFont{lmodern}{U}{euex}{m}{n}
%\DeclareMathSymbol\infty \mathord{lmodern}{"31}
\usepackage[scaled=0.875]{helvet} % ss
\usepackage{textcomp}
\usepackage{dsfont} % pour les symboles des ensembles
\usepackage[left=2cm,right=2cm,top=1.5cm,bottom=2.5cm]{geometry}
\usepackage{fancyhdr} 
\fancyhf{}
\setlength{\headheight}{26pt}
\setlength{\footskip}{18pt}
\renewcommand {\headrulewidth}{0pt}
\renewcommand {\footrulewidth}{0pt}
%\usepackage{lastpage}
%\usepackage{makeidx}
\usepackage{graphicx} % inclure le fichier graphique de pondichery 2010
\usepackage{hyperxmp}
\usepackage[colorlinks=true,pdfstartview=FitV,linkcolor=blue,citecolor=blue,urlcolor=blue]{hyperref}
\pagestyle{fancy}
\usepackage{array}
\usepackage{tabularx}
\usepackage{multirow}
\usepackage{hhline}
\usepackage[table]{xcolor}
\usepackage{arydshln}%%% lignes en tirets dans un tableau 
\usepackage{mathtools} %pour l'alignement des termes d'une matrice
\usepackage{fltpoint}
\usepackage{xcolor,colortbl}
\usepackage{pstricks,pst-all,pst-plot,pstricks-add,pst-tree,pst-3dplot,pst-func}
\usepackage{ifthen}
\usepackage{calc}
\usepackage{esvect}
\usepackage[np]{numprint}
\setlength{\parindent}{0mm}
\usepackage [alwaysadjust]{paralist}
\setdefaultenum {1.}{a)}{}{}
\hyphenpenalty=10000 
\tolerance=2000
\emergencystretch=1em
\widowpenalty=10000
\clubpenalty=10000
\raggedbottom
\renewcommand*{\tabularxcolumn}[1]{m{#1}} %centrage vertical des cellules d'un tableau tabularx
\newcommand{\Oij}{$\left( {{\mathrm{O}};\vec i,\vec j} \right)$}
\newcommand{\Oijk}{$\left( {{\mathrm{O}};\vec i,\vec j ,\vec k} \right)$}
\newcommand{\euro}{\texteuro{}}
\newcommand{\R}{\mathds {R}}
%\newcommand{\C}{\mathds {C}}
\newcommand{\N}{\mathds {N}}
\newcommand{\Z}{\mathds {Z}}
\newcommand{\Q}{\mathds {Q}}
\newcommand{\e}{\mathrm {e}}
\newcommand{\dd}{\,\mathrm{d}}
\DecimalMathComma
%%% Mise en forme des algorithmes
\usepackage[french,boxed,titlenumbered,lined,longend]{algorithm2e}
  \SetKwIF {Si}{SinonSi}{Sinon}{si}{alors}{sinon\_si}{alors}{fin~si}
 \SetKwFor{Tq}{tant\_que~}{~faire~}{fin~tant\_que}
 \SetKwFor{PourCh}{pour\_chaque }{ faire }{fin pour\_chaque}
 \SetKwInput{Sortie}{Sortie}
\newcommand{\Algocmd}[1]{\textsf{\textsc{\textbf{#1}}}}\SetKwSty{Algocmd}
  \newcommand{\AlgCommentaire}[1]{\textsl{\small  #1}}  
%%%%%%%%%%%%%%%
\newcounter{nexo}    % déclaration du numéro d'exo
\setcounter{nexo}{0} % initialisation du numero
\newcommand{\exo}{%
  \stepcounter{nexo} 
  \par{\textsf{\textbf{\textcolor{blue}{{\textsc{exercice  \arabic{nexo}}}}}}}
  }%
\newcommand{\QCM}[3]{ $\bullet$ \quad #1 \hfill 
$\bullet$ \quad #2 \hfill 
$\bullet$ \quad #3 }
\newcommand{\QCMc}[4]{   #1 \\ 
\makebox[3cm][l]{\psframe(.25,.25) \qquad #2 }\hfill 
\makebox[3cm][l]{\psframe(.25,.25) \qquad #3 }\hfill 
\makebox[3cm][l]{\psframe(.25,.25) \qquad #4 }} % QCM avec cases à cocher%
% en-tête
%\lhead{\footnotesize{\textsf{Lycée JANSON DE SAILLY\\ {\today}}}} %gauche
\chead{\textsf{\textcolor{red}{\textbf{\textsc{DS02}}}}} %centre
%\rhead{\footnotesize{\fbox{\textsf{S1E2}}}} %droit
%pied de page
%\lfoot{\scriptsize{\textsf{\textsc{A. Yallouz}}}} %gauche
\cfoot{\scriptsize{\textsf{-~\thepage{}~-}}}  %centre                                                   
%\rfoot{\htmladdnormallink{\scriptsize\textsf{\textsc{MATH@ES}}}{http://yallouz.arie.free.fr}} %droit 
\hypersetup{
pdftitle={titre du document.}, 
pdfauthor={Arié Yallouz},   
pdfsubject={Exercices Terminale ES.}, 
pdfkeywords={exercices, sujets.},
pdfcontacturl={http//yallouz.arie.free.fr/},
pdflang={fr},
baseurl={http://yallouz.arie.free.fr/serveurtex/exostex.php}
}
\begin{document}

\exo

 \medskip 

\medskip 


La fonte GS (graphite sphéroïdal) possède des caractéristiques mécaniques élevées et proche de celles des aciers. Une entreprise fabrique des pièces de fonte GS qui sont utilisées dans l'industrie automobile.

Ces pièces sont coulées dans des moules de sable et ont une température de 1400\,\degres C à la sortie du four. Elles sont entreposées dans un local dont la température ambiante est maintenue à une température de 30\,\degres C. Ces pièces peuvent être démoulées dès lors que leur température est inférieure à 650\,\degres C.

La température en degrés Celsius d'une pièce de fonte est une fonction du temps $t$, exprimé en heures, depuis sa sortie du four. On admet que cette fonction $f$, définie et dérivable sur l'intervalle $\left[ 0 ; +\infty\right[$, est une solution sur cet intervalle de l'équation différentielle :\[(E): \quad 
  y'+0,065y = 1,95.\]

Vous pourrez effectuer l'intégralité du TP avec le logiciel GEOGEBRA, EXCEL, XCAS...

Les questions où l'on vous demandera d'effectuer les calculs seront indiquées de manière explicite.

\bigskip

\begin{enumerate}
\item Résoudre l'équation homogène $(E_0)$ de cette équation différentielle

\vspace{2cm}

\item Déterminer les solutions de  l'équation différentielle $(E)$

\vspace{2cm}

\item En déduire une solution particulière $f_p$ de l'équation différentielle $(E)$

\vspace{2cm}

\item Donner la valeur de $f(0)$ et déterminer alors \textbf{par le calcul} la solution de l'équation différentielle $(E)$ qui correspond au problème donné.

\vspace{3cm}

\newpage

\item Dans cette partie, on considère la fonction $f$ définie pour tout $t \in [0;+\infty[$ par $$f(t)=1370 \e^{-0.065t}+30$$

\begin{enumerate}

\item Déterminer la fonction $f'$, dérivée de la fonction $f$.
	
	\vspace{2cm}
	
\item La pièce de fonte peut-elle être démoulée après avoir été entreposée 5 heures dans le local ?
	
\vspace{3cm}
	
	\bigskip
\end{enumerate}
	

\item
  \begin{enumerate}
  \item Déterminer au bout de combien de temps au minimum la pièce pourra être démoulée. Arrondir au centième d'heure.
  
  \vspace{3cm}
  
  
  \item Pour éviter la fragilisation de la fonte, il est préférable de ne pas démouler la pièce avant que sa température ait atteint 325\,\degres C.

    Dans ce cas faudra-t-il attendre exactement deux fois plus de temps que pour un démoulage à 650\,\degres C ? Justifier la réponse.

  \end{enumerate}
  
  \end{enumerate}
\end{document}

  


\end{enumerate}


 \bigskip 

\end{document}

 
 \end{document}